\begin{abcliste}
\abc $p = m\,v$, with the mass in kilograms (C$_{60}$: $m = \mcuO = \mc$):
\begin{align}
 p_\mathrm{e} &= \ncme\times\vel \approx \result{\pelP{0}{2}} \\
 p_\mathrm{C_{60}} &= \mc\times\vc \approx \result{\pcP{0}{2}} \\
 p_\mathrm{dust} &= \md\times\vd \approx \result{\pdP{0}{2}}
\end{align}

\abc $\lambda = h/p$:
\begin{align}
 \lambda_\mathrm{e} &= \frac{\nch}{\pel} \approx \result{\lelP{-9}{2}} \\
 \lambda_\mathrm{C_{60}} &= \frac{\nch}{\pc} \approx \result{\lcP{-12}{2}} \\
 \lambda_\mathrm{dust} &= \frac{\nch}{\pd} \approx \result{\ldP{0}{2}}
\end{align}

\abc Diffraction is visible when the wavelength is not much smaller than the spacing of the obstacles. Only the wavelength of the electron, \lelP{-9}{2}, is comparable to the spacing of the atoms: the crystal diffracts \result{\text{only the electron}}. The wavelength of C$_{60}$ is about a hundred times too short; still, with a grating of \SI{100}{\nano\meter} and a long flight path its interference has been seen (Vienna, 1999). For the grain of dust the wavelength is so absurdly short that no wave behaviour can ever be seen.
\end{abcliste}
